% ===================================================================
%  Generado por Erlen
%  Compílalo en Overleaf, o localmente con dos pasadas de pdflatex.
% ===================================================================
\documentclass[aspectratio=169,11pt]{beamer}
\usetheme{default}
% Sin barras ni adornos: el marco acompaña a las figuras y no compite con
% ellas. En vez de un acento único, tres azules emparentados: el título en el
% más oscuro, las viñetas y las reglas en el medio, los bloques en el acero.
\definecolor{acuarelaPapel}{HTML}{F7FAFB}
\definecolor{acuarelaTinta}{HTML}{1E2A31}
\definecolor{acuarelaTitulo}{HTML}{1D5C74}
\definecolor{acuarelaVineta}{HTML}{2B7B9C}
\definecolor{acuarelaBloque}{HTML}{4A6A9E}
\setbeamercolor{background canvas}{bg=acuarelaPapel}
\setbeamercolor{normal text}{fg=acuarelaTinta}
\setbeamercolor{structure}{fg=acuarelaVineta}
\setbeamercolor{frametitle}{bg=acuarelaPapel,fg=acuarelaTitulo}
\setbeamercolor{block title}{bg=acuarelaBloque!10!acuarelaPapel,fg=acuarelaBloque}
\setbeamercolor{block body}{bg=acuarelaBloque!6!acuarelaPapel}
\setbeamerfont{frametitle}{size=\large,series=\mdseries}
% La letra de pantalla es Fira Sans Light: compila con XeLaTeX o LuaLaTeX y
% añade \usepackage[light]{FiraSans} para clavarla. Con pdfLaTeX sale la sans
% de fábrica, que también vale.
% El filete de tres colores bajo el título y en la portada es de Erlen:
% Beamer no lo trae.
\usepackage[utf8]{inputenc}
\usepackage[T1]{fontenc}
% Español: se carga solo si babel-spanish está instalado (Overleaf lo trae).
\IfFileExists{spanish.ldf}{\usepackage[spanish,es-nodecimaldot]{babel}}{}
\usepackage{amsmath,amssymb}
\usepackage[version=4]{mhchem}
\definecolor{erlenacento}{HTML}{2B7B9C}
\usepackage{tikz}
\usetikzlibrary{arrows.meta, calc}
\usepackage{pgfplots}
\pgfplotsset{compat=1.18}
\definecolor{serieA}{HTML}{0072B2}
\definecolor{serieB}{HTML}{E69F00}
\definecolor{serieC}{HTML}{009E73}
\definecolor{serieD}{HTML}{D55E00}
\definecolor{serieE}{HTML}{CC79A7}
\definecolor{serieF}{HTML}{56B4E9}
\renewcommand{\figurename}{Figura}
\renewcommand{\tablename}{Tabla}
\usepackage{siunitx}
\sisetup{per-mode=symbol, range-units=single}
\setbeamertemplate{footline}[frame number]
\setbeamertemplate{navigation symbols}{}
% Algunos temas (Metropolis entre ellos) insertan solos una diapositiva al
% empezar cada sección. Erlen ya emite la suya, así que se desactiva la
% automática: de otro modo cada sección saldría dos veces y la numeración
% del PDF dejaría de coincidir con la de la app.
\AtBeginSection{}
\AtBeginSubsection{}

\title[Cinética de primer orden]{Cinética de primer orden}
\subtitle{Ejemplo didáctico · datos ilustrativos}
\date{\today}

\begin{document}

% --- diapositiva 1 ---
\begin{frame}[plain]
  \transdissolve[duration=0.25]
  \titlepage
\end{frame}
\note{%
  {\scriptsize\textbf{Tiempo previsto: 0:30 min}}\par\medskip
  Ejemplo didáctico. Sustituye contenido y datos por los de tu trabajo. No atribuir estas cifras a un experimento.\par
}

% --- diapositiva 2 ---
\begin{frame}{La pregunta que guía el trabajo}
  \transdissolve[duration=0.25]
  {\large ¿Describe un modelo de primer orden la disminución de concentración?}
  \medskip

  {\small Caso de uso ilustrativo. Define aquí el sistema, la comparación y el alcance.}
\end{frame}
\note{%
  {\scriptsize\textbf{Tiempo previsto: 1:30 min}}\par\medskip
}

% --- diapositiva 3 ---
\begin{frame}{Una estrategia explícita}
  \transdissolve[duration=0.25]
  \begin{figure}
    \centering
    \definecolor{sa0f}{HTML}{CFDCE6}
    \definecolor{sa1f}{HTML}{9DB7CC}
    \definecolor{sa2f}{HTML}{6B93B2}
    \definecolor{sa3f}{HTML}{396E98}
    \definecolor{sat1843819360}{HTML}{15181C}
    \definecolor{sat1226267613}{HTML}{FFFFFF}
    % Con babel-spanish las palabras largas se parten solas y el texto respira mejor.
    \begin{tikzpicture}[x=1cm,y=1cm,font=\footnotesize]
      \draw[fill=sa0f, draw=none] (0.000,0.000) -- (2.226,0.000) -- (2.518,-1.254) -- (2.226,-2.508) -- (0.000,-2.508) -- cycle;
      \draw[fill=sa1f, draw=none] (2.644,0.000) -- (4.870,0.000) -- (5.162,-1.254) -- (4.870,-2.508) -- (2.644,-2.508) -- (2.936,-1.254) -- cycle;
      \draw[fill=sa2f, draw=none] (5.288,0.000) -- (7.514,0.000) -- (7.806,-1.254) -- (7.514,-2.508) -- (5.288,-2.508) -- (5.580,-1.254) -- cycle;
      \draw[fill=sa3f, draw=none] (7.932,0.000) -- (10.157,0.000) -- (10.450,-1.254) -- (10.157,-2.508) -- (7.932,-2.508) -- (8.224,-1.254) -- cycle;
      \node[text=sat1843819360, align=center, text width=1.75cm, inner sep=1pt] at (1.113,-1.254) {\scriptsize \bfseries Registrar tiempo y concentración};
      \node[text=sat1843819360, align=center, text width=1.75cm, inner sep=1pt] at (3.976,-1.254) {\bfseries Declarar el modelo};
      \node[text=sat1843819360, align=center, text width=1.75cm, inner sep=1pt] at (6.620,-1.254) {\bfseries Estimar k};
      \node[text=sat1226267613, align=center, text width=1.75cm, inner sep=1pt] at (9.264,-1.254) {\bfseries Revisar residuos y controles};
    \end{tikzpicture}
  \end{figure}
\end{frame}
\note{%
  {\scriptsize\textbf{Tiempo previsto: 1:30 min}}\par\medskip
  Explica las decisiones y los controles, no solo los nombres de las técnicas.\par
}

% --- diapositiva 4 ---
\begin{frame}{Mostrar la evidencia}
  \transdissolve[duration=0.25]
  \begin{figure}
    \centering
    \begin{tikzpicture}
      \begin{axis}[
        width=0.90\linewidth,
        height=0.50\linewidth,
        xlabel={Tiempo (min)},
        ylabel={c/c₀},
        grid=major,
        grid style={line width=.2pt, draw=gray!25},
        tick align=outside,
        tick pos=left,
        axis line style={gray!60},
        label style={font=\small},
        tick label style={font=\footnotesize}
      ]
        \addplot[serieA, line width=1pt, mark=none, smooth] coordinates {
        (0,1) (2,0.67) (4,0.449) (6,0.301) (8,0.202) 
        };
      \end{axis}
    \end{tikzpicture}
    \caption{Datos ilustrativos para explicar el método; sustituir por mediciones propias.}
  \end{figure}
  \medskip

  \begin{equation*}
    c(t)=c_0 e^{-kt},\quad k=0.2\;\mathrm{min}^{-1}
  \end{equation*}
\end{frame}
\note{%
  {\scriptsize\textbf{Tiempo previsto: 1:30 min}}\par\medskip
  Datos o modelos didácticos: reemplazar antes de presentar resultados propios.\par
}

% --- diapositiva 5 ---
\begin{frame}{Interpretación y límites}
  \transdissolve[duration=0.25]
  \begin{itemize}
    \item Valores ilustrativos calculados con una exponencial y redondeados.
    \item La curva no identifica por sí misma un mecanismo químico.
    \item Comparar modelos alternativos y residuos antes de concluir.
  \end{itemize}
\end{frame}
\note{%
  {\scriptsize\textbf{Tiempo previsto: 1:30 min}}\par\medskip
}

% --- diapositiva 6 ---
\begin{frame}{Antes de compartir}
  \transdissolve[duration=0.25]
  \begin{itemize}
    \item Origen de t = 0 y resolución temporal
    \item Temperatura, unidades de k y réplicas
    \item Criterio de ajuste e incertidumbre de parámetros
  \end{itemize}
  \medskip

  {\small Conserva el proyecto JSON y revisa la exportación final.}
\end{frame}
\note{%
  {\scriptsize\textbf{Tiempo previsto: 1:30 min}}\par\medskip
}

\end{document}